\documentclass[11pt]{article}

\usepackage[margin=1in]{geometry}
\usepackage{graphicx}
\usepackage{booktabs}
\usepackage{amsmath}
\usepackage[hidelinks]{hyperref}
\usepackage[authoryear,round]{natbib}
\emergencystretch=3em

\title{Reliability modeling for transportable perturbation-response predictions in public single-cell functional-genomics screens}
\author{Yiheng Yang\\
School of Future Technology, Dalian University of Technology\\
School of Bioengineering, Dalian University of Technology\\
\texttt{yangyiheng@mail.dlut.edu.cn}}
\date{}

\begin{document}
\maketitle

\begin{abstract}
Computational models are increasingly used to infer gene-perturbation responses from single-cell functional genomics screens, yet their predictions can fail when transferred to unseen perturbations, low-support contexts, or external datasets. Such failures can mislead downstream pathway interpretation and perturbation prioritization. We present the Perturbation Transportability Layer (PTL), a model-agnostic reliability model for already-generated single-cell perturbation-response outputs. PTL combines reference-aware delta signatures, context-stress splits, relative transportability labels, deployment-feature keep/abstain modeling, and a failure-mode atlas. We evaluated PTL on public Perturb-seq resources, including an independent public GSE284197 screen and additional public-screen holdouts. Across 525 transparent baseline runs, random-split rankings did not transport across held-out perturbation, combination, low-support, dataset-heldout, and external-holdout regimes. In the primary setting, the full random-forest PTL model reduced the mean false-transportability rate from 0.71 for naive confidence to 0.24. Retained predictions showed higher transportability and perturbation-retrieval diagnostics than rejected high-risk outputs, while pathway-level summaries remained diagnostic rather than uniformly improved. Case examples showed that high-confidence predictions can still reverse or degrade top-gene signals under perturbation and dataset transfer. A signature-level GEARS adapter provided an output-contract compatibility demonstration, not a performance benchmark. PTL provides a reliability layer evaluated across public-screen holdouts for prioritizing perturbation predictions before downstream biological interpretation.
\end{abstract}

\noindent\textbf{Keywords:} functional genomics; Perturb-seq; single-cell perturbation prediction; transportability; reliability modeling; dataset shift

\section{Introduction}

Single-cell functional genomics screens now make it possible to connect genetic perturbations with transcriptional response phenotypes at scale \citep{dixit2016perturbseq,adamson2016parseq,datlinger2017pooled,xie2017multiplexed,norman2019manifolds,srivatsan2020sciplex,datlinger2021ultra,papalexi2021eccite,frangieh2021multimodal,replogle2022genome,peidli2024scperturb}. These Perturb-seq and related CRISPR screens are increasingly used to study gene function, pathway regulation, combinatorial genetic interactions, and candidate target prioritization, and curated perturbation resources are growing rapidly \citep{perturbase2025,perturbdb2025}. At the same time, computational perturbation-response predictors such as scGen, CPA, GEARS, CellOT, CINEMA-OT, and related single-cell models have made it feasible to infer responses for perturbations or contexts that are only partially observed in the training data \citep{lotfollahi2019scgen,lotfollahi2023cpa,roohani2023gears,bunne2023cellot,dong2023cinemaot}. Broader single-cell representation models and virtual-cell proposals further increase interest in prediction systems that can be interrogated before biological interpretation \citep{lopez2018scvi,gayoso2022scvitools,theodoris2023geneformer,cui2024scgpt,bunne2024virtualcell}. For biological interpretation, however, the key question is not only whether a model has good average accuracy, but whether a specific prediction remains reliable when it is used to prioritize a pathway, perturbation, or follow-up experiment.

Random-split accuracy can obscure this reliability problem. In a random split, train and test data may share perturbations, control strata, cell contexts, support structure, and experimental protocols. Recent perturbation-prediction benchmarks and subtask-decomposition work similarly emphasize that models can behave differently across biologically distinct generalization regimes \citep{generalizable2025benchmark,stamp2024}. Dataset shift, domain generalization, calibration, uncertainty, selective prediction, and conformal prediction literatures all emphasize that average in-distribution accuracy does not guarantee reliable deployment under transfer \citep{zhou2022dgsurvey,koh2021wilds,ovadia2019uncertainty,guo2017calibration,hendrycks2016ood,geifman2017selective,vovk2005algorithmic,angelopoulos2021conformal}. A model can therefore appear strong while failing on the settings where computational prediction is most useful: unseen perturbations, unseen combinations, low-support signatures, new datasets, or external public screens. These failures are biologically consequential because a non-transportable prediction can enter downstream pathway enrichment, perturbation retrieval, or gene-function interpretation as if it were an experimentally supported response. Reliability modeling for perturbation prediction should therefore ask where transfer fails and whether high-risk predictions can be filtered before biological interpretation.

We address this problem with the Perturbation Transportability Layer (PTL), a model-agnostic reliability model for already-generated single-cell perturbation-response predictions. PTL does not replace an expression predictor; it estimates the risk that an already-generated response prediction will fail under transfer. PTL uses four linked components: a stress-split benchmark for perturbation, combination, support, dataset, and external transfer; a relative transportability labeling rule anchored to split-specific difficulty; a deployment-feature keep/abstain model; and a failure-mode atlas that links prediction errors to novelty, support, and dataset boundaries. Using public Perturb-seq screens, we show that random-split winners do not define universal transfer winners, that PTL reduces false transportability compared with confidence filtering, and that severe failures concentrate at biologically interpretable transfer boundaries.

\section{Methods}

\subsection{Overview of PTL reliability modeling}

PTL is a reliability model for already-generated single-cell perturbation predictions. Its intended input is a perturbation-response dataset or benchmark manifest containing an expression matrix, cell metadata, perturbation labels, control labels, dataset identifiers, context annotations, support counts, and prediction outputs from one or more models. Its output is a set of context-stress benchmark scores, per-signature transportability labels, a model-agnostic keep/abstain score, and a failure-mode atlas.

The analysis has six modules: input standardization, reference-aware signature construction, context-stress split generation, relative transportability labeling, deployment-feature reliability modeling, and failure-mode reporting. This organization keeps the expression predictor separate from the reliability model and makes each audit artifact traceable.

\subsection{Step-by-step PTL protocol}

The protocol is intended to be run as an auditable sequence rather than as a single opaque benchmark command. It prepares a perturbation-response object, verifies matched controls, constructs delta signatures, writes split manifests and leakage audits, runs a predictor or adapter, assigns relative labels, trains PTL, and reports risk--coverage, calibration, failure-atlas, and biological diagnostic tables. A run should not be interpreted until the split audit, gene-space audit, feature audit, and output-contract checks have passed.

\begin{center}
\fbox{\begin{minipage}{0.94\linewidth}
\textbf{Algorithm 1. PTL input-output contract and training procedure}

\textbf{Input:} expression matrix, cell metadata, perturbation/control labels, dataset/context labels, split manifest, and model prediction outputs.

\textbf{Output:} aligned signatures, transportability labels, keep/abstain scores, risk--coverage summaries, and a failure-mode atlas.

\begin{enumerate}
\item Align prediction and target signatures in a shared gene space.
\item Construct reference-aware delta signatures, \(\Delta_{p,c}=\bar{x}_{p,c}-\bar{x}_{0,c}\).
\item Generate context-stress splits and audit support, leakage, and train--test distance.
\item Assign relative labels with \(I[\cos(\hat{\Delta}_{p,c},\Delta_{p,c}) \ge \tau \times m_s]\), where \(\tau\in\{0.6,0.7,0.8,0.9\}\).
\item Train the keep/abstain reliability model using deployment-available features only.
\item Report false-transportability, risk--coverage, calibration, and failure-atlas outputs.
\end{enumerate}
\end{minipage}}
\end{center}

Here, \(m_s\) denotes the median random-split anchor fidelity computed within the corresponding dataset, model, and seed context before applying the stress-family label rule.

\subsection{Input contract and data resources}

The internal benchmark surface used harmonized public scPerturb resources, including NormanWeissman2019 and ReplogleWeissman2022 screens \citep{norman2019manifolds,replogle2022genome,peidli2024scperturb}. The independent public experimental validation source was the GSE284197 Perturb-seq screen \citep{geo2025gse284197}, processed with the same metadata and signature contract but held out from the internal benchmark surface. This manuscript treats GSE284197 as public experimental validation, not as newly generated experimental data. Standard single-cell preprocessing choices were aligned with common Scanpy, Seurat, and scVI/scvi-tools practice where relevant \citep{wolf2018scanpy,stuart2019integration,hafemeister2019sctransform,lopez2018scvi,gayoso2022scvitools}.

Each dataset was required to expose expression values, feature names, perturbation or control labels, dataset identity, and enough context metadata to identify matched control strata. Prediction outputs were required to provide a test prediction matrix, a matched target matrix, test metadata, a gene list, and run-level metadata. This contract allowed transparent baselines and an adapted published architecture to be evaluated through the same downstream PTL machinery.

\subsection{Reproducible preprocessing and benchmark parameters}

All public screens were processed into a common signature contract before modeling. Input objects were required to contain a normalized expression matrix, gene symbols, perturbation labels, control labels, dataset identifiers, and context fields sufficient to match perturbed cells to controls. Gene filtering used the 6,257-gene global intersection used in this benchmark. Cells were aggregated into perturbation--context pseudobulk signatures only when the corresponding perturbed and matched-control strata were available; signatures without an eligible reference were excluded from predictor scoring and recorded during the split audit. Low-support stress was defined from signature-level backing cell counts: non-control signatures were sorted by \(n_{\mathrm{cells}}\), and group keys with \(n_{\mathrm{cells}}\) at or below the empirical 20th-percentile cutoff were assigned to the low-support test mask and then projected onto the cell-level track.

Splits were generated deterministically from manifest files and seeds 0--2 where available. The random split served as the anchor; stress splits held out perturbations, perturbation combinations, low-support signatures, complete datasets, or the independent public screen. Each manifest recorded train/test membership, held-out units, support counts, perturbation overlap, reference-control overlap, dataset overlap, and train--test context distance. Baseline predictors used fixed hyperparameters across datasets: ridge regression used the same regularization grid selected within the training split, k-nearest-neighbor transfer used context-distance neighbors, and mean-delta baselines used training-only perturbation or global deltas. Model-confidence features were restricted to target-free, deployment-available quantities such as support, distance-to-training context, predicted-effect norm, and model-internal or proxy uncertainty where available. Outcome-derived fidelity metrics were excluded from PTL features and used only for evaluation.

Biological diagnostics were computed from the same aligned prediction and target signatures. Pathway summaries used the project gene-set collection recorded in \path{transportability_gene_sets.json}; perturbation retrieval compared each prediction with held-out target signatures in the same evaluation surface; false transportability was the fraction of accepted predictions whose relative transportability label was negative. Calibration curves, ROC/precision--recall summaries, and risk--coverage curves were computed from out-of-fold PTL scores so that selective-prediction evidence did not reuse the same labels for fitting and reporting.

\subsection{Signature construction}

The canonical modeling surface is the signature track. Each non-control signature represents a perturbation and context stratum with a matched reference control stratum. Let \(x_{i}\) denote an expression vector for cell \(i\), \(p\) a perturbation label, and \(c\) a context stratum. The perturbation signature is the pseudobulk mean \(\bar{x}_{p,c}\), and the matched control signature is \(\bar{x}_{0,c}\). The delta signature is
\[
\Delta_{p,c}=\bar{x}_{p,c}-\bar{x}_{0,c}.
\]
Predictions and targets were compared on these delta signatures in a shared gene space. Cell-level information was retained for support counts, leakage checks, and future cell-level extensions, but the claims here are restricted to signature-level prediction.

\subsection{Context-stress split generator}

PTL requires an explicit split manifest before model training. The random split serves as a low-stress anchor. Stress splits hold out perturbations, perturbation combinations, low-support signatures, whole datasets, or an independent public experimental dataset. In algorithmic terms, the split generator selects a stress family, identifies the corresponding withheld units, writes deterministic train/test membership across seeds, audits leakage of perturbation and context labels, and records support and context-distance summaries.

The expanded benchmark surface used the 6,257-gene global intersection used in this benchmark for transparent baseline runs. The completed manifest contained 525 transparent baseline rows across five baseline models, three seeds where available, and the six split families: random, held-out perturbation, held-out combination, low-support, dataset-heldout, and external-holdout.

\subsection{Baseline and adapter predictors}

Five transparent full-gene baselines were evaluated: a control-mean baseline, a global-delta baseline, a perturbation-mean-delta baseline, a cell-context k-nearest-neighbor delta baseline, and ridge regression. These baselines were chosen because they expose reference behavior, average perturbation effects, context-aware nearest-neighbor transfer, and linear multivariate learning without hiding the reliability analysis behind a specialized architecture.

To test whether the PTL contract can accept a published perturbation predictor, we also implemented a GEARS adapter \citep{roohani2023gears}. The adapter converts signature-level AnnData objects into the GEARS data interface, uses a custom train/test split, trains GEARS, and writes the same prediction matrix, metadata, gene-list, log, and registry contract as the transparent baselines. Because this adapter demonstration uses signature-level pseudobulk profiles and conservative graph fallbacks, it is reported as an output-contract compatibility demonstration, not a performance benchmark.

\subsection{Transportability labeling}

The primary fidelity metric was mean cosine similarity on non-control test signatures. Per-signature transportability labels were defined with a relative anchor rule: a prediction was labeled transportable when its cosine similarity was at least \(0.8\) times the median random-split anchor fidelity computed for the matching dataset, model, and seed context. This relative rule avoids imposing a single absolute threshold across regimes with very different intrinsic difficulty. Threshold sensitivity and risk--coverage summaries were retained as robustness diagnostics because the label is intended to support operational filtering, not mechanistic truth.

\subsection{PTL training and deployment}

PTL is trained after a predictor has produced outputs. It does not predict expression. It estimates whether an already-generated prediction should be retained. Deployment-available feature groups included model and split identifiers, dataset identifiers, target-free model-confidence proxies, support counts, perturbation novelty, combination novelty, reference availability, predicted effect-size summaries, stress-family metadata, and optional train-test context-distance summaries. A feature audit excluded target, label, and post-outcome columns before training. Identifier features were interpreted cautiously and accompanied by ablations to test whether PTL learned generalizable reliability structure rather than only memorizing dataset or split labels.

Logistic regression was used as a calibrated baseline estimator and random forest as the nonlinear comparator. Ablations removed context-distance, perturbation-novelty, uncertainty, or optional pathway features. The primary reliability outcome was false-transportability rate among accepted predictions; ROC AUC, average precision, calibration, risk--coverage, and selective-risk summaries were reported as supporting diagnostics.

\subsection{Biological summaries and failure-mode atlas}

Biological usefulness was assessed with public-data summaries rather than newly generated experiments. Gene-set summaries measured pathway-level cosine and direction consistency using the available pathway collection. Perturbation-retrieval summaries tested whether predicted signatures preserved perturbation identity relative to other test signatures. These metrics are not substitutes for gene-level fidelity, but they provide a biological check on retained versus rejected outputs.

The failure-mode atlas groups per-signature examples by model, split family, support and novelty context, transportability status, risk, and severity. The atlas is descriptive: it identifies recurring regimes where predictions degrade, especially severe failure, high-risk transportability errors, novelty-associated degradation, and low-support degradation.

\subsection{Comparison with alternative reliability approaches}

PTL is not the only possible way to filter perturbation predictions. Table~\ref{tab:methodcomparison} summarizes the operational differences between common alternatives and the proposed reliability approach. Random-split benchmarking and external-only validation are useful evaluation views but do not provide per-signature keep/abstain decisions. Confidence filtering and support/novelty heuristics can abstain, but they do not jointly use split semantics, model outputs, and deployment features. PTL combines these inputs into a model-agnostic reliability layer and a failure-mode report.

\begin{table}[t]
\centering
\caption{Operational comparison between PTL and alternative reliability approaches.}
\label{tab:methodcomparison}
\resizebox{\linewidth}{!}{%
\input{../results/tables/methods_comparison_table.tex}
}
\end{table}

\subsection{Practical guidance and troubleshooting}

Recommended use starts with the split audit, not with model scores. Users should verify that held-out perturbations, combinations, datasets, and controls match the intended biological question; confirm that gene spaces are aligned; inspect support distributions; and run simple baselines before interpreting a complex model. If context-distance features add little predictive value, as in the present benchmark, they should be retained as diagnostic summaries but not over-interpreted as necessary PTL features. If a published model cannot ingest the full expression contract, an adapter may still be useful, but claims should be restricted to the surface actually tested.

\section{Results}

\subsection{Benchmark design for perturbation-response transportability}

We organized PTL as a computational reliability study for public single-cell functional genomics screens (Figure~\ref{fig:concept}). The analysis starts from perturbation-response signatures and model outputs, evaluates them under stress-transfer regimes, and assigns keep/abstain reliability scores before biological interpretation. The benchmark surface contained 525 transparent baseline runs with consistent gene counts and valid metric outputs. Random splits were retained as low-stress anchors, whereas held-out perturbation, held-out combination, low-support, dataset-heldout, and external-holdout splits tested biologically relevant transfer questions.

Figure~\ref{fig:benchmark} summarizes the dataset surface and split audit. The independent public GSE284197 screen was treated as an external experimental validation source, and leave-one-screen-out dataset-heldout summaries were used to test additional public experimental screens. The hardest stress family by average non-control cosine was dataset-heldout transfer, with mean cosine 0.049 across implemented baselines. GSE284197 showed high composite difficulty because perturbation overlap, dataset identity, and train-test context all shifted simultaneously. These audits define the biological setting of the study: predicted signatures are not only evaluated for average fidelity, but for whether fidelity survives perturbation and dataset transfer.

\begin{figure}[t]
\centering
\includegraphics[width=\linewidth]{../results/figures_cbac/fig1_simplified_framework.png}
\caption{PTL framework for reliability modeling in public single-cell perturbation screens. Public Perturb-seq inputs are converted into reference-aware delta signatures, evaluated with context-stress split families, passed through a model-agnostic PTL reliability model, and reported as keep/abstain decisions for downstream biological use.}
\label{fig:concept}
\end{figure}

\begin{figure}[t]
\centering
\includegraphics[width=\linewidth]{../results/figures_cbac/fig2_cbac_benchmark_design.png}
\caption{Benchmark composition, split construction, and compact audit diagnostics. The panels report the public screen surface, split-family transfer questions, and overlap/leakage checks before model performance is interpreted. Support and context-distance details are provided in Supplementary Figure S1.}
\label{fig:benchmark}
\end{figure}

\subsection{Model rankings collapse under perturbation and dataset transfer}

Although ridge regression led the random split and two stress regimes, random-split ranking did not define a universal winner across transfer regimes (Table~\ref{tab:familywinners}, Figure~\ref{fig:decay}). Ridge regression led the random-split, low-support-split, and unseen-perturbation-split conditions. In contrast, global-delta baseline led the dataset-heldout split and external holdout, whereas perturbation-mean-delta led the unseen-combination split. Corrected paired comparisons found 17 FDR-significant model contrasts, supporting the conclusion that the rank shifts were not only descriptive artifacts of one summary table.

\begin{table}[t]
\centering
\caption{Family-level winners from the transfer-decay summary.}
\label{tab:familywinners}
\begin{tabular}{lll}
\toprule
Split family & Leading model & Mean non-control cosine \\
\midrule
random\_split & ridge\_regression\_baseline & 0.373 \\
low\_support\_split & ridge\_regression\_baseline & 0.307 \\
unseen\_perturbation\_split & ridge\_regression\_baseline & 0.311 \\
unseen\_combination\_split & perturbation\_mean\_delta\_baseline & 0.664 \\
dataset\_heldout\_split & global\_delta\_baseline & 0.056 \\
external\_holdout & global\_delta\_baseline & 0.112 \\
\bottomrule
\end{tabular}
\end{table}

\begin{figure}[t]
\centering
\includegraphics[width=\linewidth]{../results/figures_cbac/fig3_cbac_ranking_instability.png}
\caption{Model ranking instability across transfer regimes. Mean non-control cosine, rank badges, rank-shift trajectories, and dataset/external transfer behavior show that random-split ranking does not define a universal perturbation-response predictor.}
\label{fig:decay}
\end{figure}

External and public-screen holdouts reinforced the same point. On GSE284197, global-delta, cell-context k-nearest-neighbor, and perturbation-mean-delta baselines reached mean non-control cosine around 0.11, whereas ridge regression fell below zero. Additional public-screen holdouts also varied substantially: mean non-control cosine was 0.038 for ReplogleWeissman2022 RPE1, -0.025 for PapalexiSatija2021 ECCITE-seq RNA, and 0.006 for DatlingerBock2021, with the best model changing across screens. Thus, public functional genomics screens expose transfer behavior that random-split reporting compresses away.

\subsection{PTL filters false transportability better than confidence}

PTL attaches a reliability decision to predicted signatures rather than replacing the perturbation predictor (Figure~\ref{fig:ptl}). The primary deployment model was the full PTL random forest, which reduced mean false-transportability rate from 0.71 for naive confidence to 0.24, with ROC AUC 0.93 and average precision 0.93. The best-performing ablation was the no-context-distance random forest, with a similar false-transportability rate of 0.23 and ROC AUC 0.94; we treat that row as a robustness check rather than the primary claim.

\begin{figure}[t]
\centering
\includegraphics[width=\linewidth]{../results/figures_cbac/fig4_cbac_reliability_filtering.png}
\caption{PTL reliability filtering of perturbation-response predictions. The panels emphasize the full PTL random-forest reduction in false transportability, ROC and precision--recall behavior, risk--coverage, and selected reliability baselines.}
\label{fig:ptl}
\end{figure}

The selective-prediction result should be read as a transportability filter. At the default \(0.8\times\) split-specific anchor threshold, the full PTL random forest achieved a 67\% relative reduction in mean false transportability. In the risk--coverage summary, the same model accepted 80\% of predictions at the reported high-coverage operating point with false-transportability rate 0.47, illustrating the trade-off between coverage and selective risk. A calibrated logistic baseline was the next strongest comparator among the reliability baselines, with mean false-transportability rate 0.27 and ROC AUC 0.87, whereas split-family-only, support-only, novelty-only, and support-plus-novelty models remained weaker. Threshold sensitivity was stable across anchor multipliers 0.6, 0.7, 0.8, and 0.9 (Supplementary Figure S2), supporting the relative anchor rule as an operational filter rather than a fragile cutoff.

\subsection{Failure-mode atlas reveals biological transfer boundaries}

The failure-mode atlas connects transportability errors to interpretable transfer contexts (Figure~\ref{fig:failure}). Severe failures concentrated around perturbation novelty, low support, and dataset boundaries rather than appearing uniformly across the benchmark. The most extreme row was ridge regression under external holdout, where severe failures accounted for 73.4\% of signatures and perturbation overlap was near zero. Dataset-heldout severe failures were also dominated by perturbation-novelty signals. These patterns suggest that perturbation prediction reports should include failure composition and context enrichment, not only aggregate cosine.

Case examples make the biological interpretation risk concrete. Cases were selected from predefined retained, high-confidence failure, and dataset-boundary failure categories using the transportability label, PTL risk decision, and severity class. A retained PapalexiSatija2021 JAK2 perturbation prediction had cosine 0.90 and preserved all top-gene directions among the eight strongest true-response genes, sharing interferon-response genes including CD74, GBP1, GBP5, HLA-DRA, IRF1, TAP1, and WARS. In contrast, a high-confidence NormanWeissman2019 MAML2 prediction was non-transportable despite confidence 1.00, with cosine -0.22 and only two shared top genes. A dataset-heldout PapalexiSatija2021 SMAD4 case degraded further, with cosine -0.45 and top-gene direction consistency 0.25. These examples show why high-confidence predictions should not automatically enter pathway interpretation or perturbation prioritization.

\begin{figure}[t]
\centering
\includegraphics[width=\linewidth]{../results/figures_cbac/fig5_cbac_failure_cases.png}
\caption{Failure-mode atlas and biological case examples. The panels separate the failure taxonomy, model--split severity, severity composition, context enrichment, and case examples of retained and rejected perturbation predictions.}
\label{fig:failure}
\end{figure}

Retained predictions had a higher transportable-label rate than rejected high-risk predictions (0.91 versus 0.17) and stronger perturbation-identity retrieval (0.086 versus 0.034 top-1 retrieval). Pathway-level cosine was not uniformly higher for retained predictions, so we treat it as a diagnostic summary rather than a one-directional success metric. The atlas is therefore a biological audit surface: it indicates which predictions are plausible candidates for interpretation and which ones should be flagged as transfer-boundary failures.

\subsection{Adapter demonstration and generality}

The GEARS adapter produced valid PTL-compatible outputs for completed Replogle demonstration rows, including random-split rows and a 10-epoch held-out-perturbation RPE1 row (Supplementary Figure S3). These rows wrote the required prediction matrix, target matrix, metadata, gene list, run log, and metrics contract. This section tests output-contract compatibility across random and held-out perturbation settings, not GEARS competitiveness or final GEARS performance.

We also audited GSE264667/NadigOConner2024 HepG2 and Jurkat Perturb-seq as candidate external resources. The GEO h5ad files are publicly available but large, and local scPerturb proxy files expose perturbation labels, control cells, gene matrices, and metadata sufficient for contract inspection. Baseline prediction rows from the raw GEO source were not available for this analysis, so these screens remain contract-passed candidates rather than performance evidence.

\section{Discussion}

This study frames single-cell perturbation prediction as a reliability modeling problem in public functional genomics data. A random split remains useful as a low-stress anchor, but it does not answer the same biological question as predicting unseen perturbations, low-support signatures, new datasets, or an independent public screen. The observed rank reversals show that a random-split winner should not be assumed to be the safest model for perturbation prioritization or pathway interpretation.

PTL addresses a different task from the perturbation predictor itself. A predictor estimates a response signature; PTL estimates whether that completed response is transportable enough to retain. This separation is useful because reliability decisions can be layered over transparent baselines and published model adapters without changing the expression predictor. In the present public-data benchmark, PTL substantially reduced false transportability relative to naive confidence and simpler deployment-feature baselines, indicating that support, novelty, split context, and model-output summaries contain actionable reliability information.

The failure-mode atlas adds biological interpretability to the reliability model. Severe failures were concentrated at novelty, low-support, and dataset-transfer boundaries, and case examples showed that high-confidence predictions can still reverse or degrade top-gene response structure. These patterns argue for reporting failure composition together with aggregate fidelity: a prediction system that improves average cosine in familiar settings may still be risky in exactly the regimes where computational inference is most attractive.

For practical reporting, perturbation-response studies should not rely on a single random-split score. We recommend reporting stress-split performance, false-transportability rate, risk--coverage behavior, and failure composition whenever predictions are used for pathway interpretation, perturbation retrieval, or prioritization. This makes reliability evidence visible before downstream biological claims are made.

PTL should be used when evaluating perturbation-response predictors under new perturbations, prioritizing predicted responses for pathway or gene-function interpretation, comparing models across split regimes, or deciding when to abstain from high-risk predictions. It is not a replacement for experimental confirmation. Rather, it is a public-data grounded reliability layer that helps decide which predicted signatures deserve downstream biological attention and which should be treated as transfer-boundary failures.

\section{Limitations}

The analysis is signature-level. It evaluates predicted delta responses after aggregation to matched perturbation/control signatures and does not model the full distribution of single-cell states within each perturbation. This design makes cross-dataset benchmarking and output-contract checks tractable, but it can miss cell-level heterogeneity and rare-state effects.

The validation uses public experimental perturbation screens rather than newly generated experimental data. GSE284197 and the additional public-screen holdouts strengthen the experimental grounding of the reliability analysis, but they are not a census of all biological, technical, or platform shifts. GSE264667/NadigOConner2024 HepG2 and Jurkat screens passed the local contract audit as candidates \citep{geo2024gse264667}, but they are not used as validation evidence until raw-source processing and baseline rows are available.

The baseline set is transparent by design. It includes simple full-gene predictors and ridge regression, but it does not constitute a complete ranking of the published perturbation-model literature. The GEARS implementation is an adapter demonstration on the signature contract, with completed random and held-out perturbation rows, limited CPU runs, and conservative graph fallbacks, rather than a tuned GEARS model-comparison study.

PTL also depends on labeled benchmark outputs for calibration. A new disease setting, species, platform, or perturbation technology may require recalibration and renewed stress-split validation before the same operating threshold is used for biological interpretation.

Finally, the failure-mode atlas is descriptive. It identifies where predictions degrade and whether degradation is associated with novelty, support, or dataset boundaries, but it does not by itself explain the molecular mechanism of each failure. PTL should therefore be used to guide reliability-aware evaluation, prioritization, and follow-up study design.


\section*{Data and code availability}
The minimal code release for peer review is available as a tagged source snapshot at \url{https://github.com/Neabigmo/PTL/tree/v1.0-cbac-submission}. It contains the split-generation code, predictor output contract, GEARS adapter code, PTL training and evaluation code, deterministic figure scripts, environment notes, and expected-output manifests. For review, the submission package includes a complete processed evidence archive with split manifests, baseline metrics, PTL reliability outputs, case-example tables, figure source data, deterministic figure scripts, checksums, and repository metadata. Raw public datasets are not redistributed and remain available from their original repositories, including the scPerturb Zenodo release, the independent public GSE284197 Perturb-seq screen, and the GSE264667/NadigOConner2024 candidate screens listed in the external-candidate audit.

\section*{Acknowledgements}
The author thanks the maintainers and contributors of the public single-cell perturbation data resources used in this study.

\section*{Funding}
This research did not receive any specific grant from funding agencies in the public, commercial, or not-for-profit sectors.

\section*{Declaration of competing interest}
The author declares no competing interests.

\section*{CRediT author statement}
Yiheng Yang: Conceptualization, Methodology, Software, Validation, Formal analysis, Data curation, Writing -- original draft, Writing -- review and editing, Visualization.

\section*{Ethics statement}
This study reanalyzes public, de-identified single-cell perturbation datasets and does not report newly generated human or animal experimental data.

\section*{Sex and gender statement}
Sex- and gender-specific analyses were not performed because the study operates on aggregated public perturbation-response signatures and the required individual-level sex or gender annotations were not consistently available across the benchmark sources.

\section*{Declaration of generative AI and AI-assisted technologies in the manuscript preparation process}
During the preparation of this work the author used OpenAI Codex to assist with drafting, code organization, LaTeX troubleshooting, and language polishing. No generative AI or AI-assisted image-generation tools were used to create or modify the submitted figures or graphical abstract. All submitted figures and the graphical abstract were created from deterministic plotting scripts and vector elements. After using Codex, the author reviewed and edited the content as needed and takes full responsibility for the content of the manuscript.

\bibliographystyle{plainnat}
\bibliography{references}

\end{document}
